% Appendix: proofs for the framework section (_d_framework.tex). Source: docs/formalization.md, Sections 5.3, 6.2-6.4,
% and 7 (frozen). The EMBER checks and caveats belong to the instances and evaluation sections, not here.

This appendix proves the results of \Cref{sec:framework}. Throughout, an edit takes $z$ to $z'$ with written set
$W$ and item delta $\Delta = (I^-, I^+, I^\sim, I^\to)$.

\paragraph{Finding the changed items (Theorem~\ref{thm:fw:side-effects}, part i).}
Let the edit write no locating position. By (R3), every parsed item of $z$ survives in $z'$ with the same identity,
members, and locating positions, and no parsed item appears, except that an end-anchored item's members gain or lose
positions at the end, which lie in $W$. A parsed item whose members avoid $W$ is unchanged, since $z'$ equals $z$
outside $W$. A parsed item whose members meet $W$ takes its new content from $z'$ and belongs to $I^\sim$ if that
content differs. For each scanned kind, parts (i) and (ii) of (R4) say that the items outside the margin are the same
in $z$ and $z'$. The items of $z$ that overlap the margin are read from the item table, those of $z'$ are found by the
recognizer on the margin alone (part iii), and comparing the two by identity gives the kind's part of $\Delta$. For a
declared structural edit, (R5) says that the declared delta is the parsed part, and the scanned kinds are handled as
above with the declared range in place of the margin.

\paragraph{Adjusting the sums (Theorem~\ref{thm:fw:side-effects}, parts ii--iv).}
(ii) The items of $z'$ split into the unchanged items, $I^+$, the modified items after the edit, and the moved items;
those of $z$ into the unchanged items, $I^-$, the modified items before the edit, and the moved items. An additive sum
splits over any partition of the items it adds. The unchanged items have the same content in both inputs, and so, by
(R2), the same contributions; a moved item keeps its content, hence its contribution. Subtracting the two splits gives
the formula. (iii) A pool's value over $z'$ is the maximum over the unchanged items, which the item table lists, and
the added and modified ones, which $\Delta$ gives. The old value alone does not suffice: pools holding $\{5, 3\}$ and
$\{5, 5\}$ both have value 5, and removing an item of value 5 leaves 3 in the first and 5 in the second. (iv) By (R1).
Applying (i)--(iv) once per edit, with the item table carried forward, gives the side effects of a sequence of edits.

\paragraph{Signatures of operations stated in elements.}
Let an operation stated in elements, on every input of a population $F$ and for every parameter, be in place and
write no locating position. By (R3), no parsed item is added, removed, or moved, and a parsed item whose members avoid
$W$ keeps its content; by (R4), no scanned item changes outside the margins. The operation's signature on an input is
therefore contained in the modifications of parsed items whose members $W$ meets, and the changes of scanned items
that overlap the margins, both read from the input's item table. Its signature over $F$ is the union over the inputs.
The condition must hold on every input of $F$: on a PE file whose section table points past its end, appending writes
a locating position, and such inputs must be excluded by the precondition or treated as structural.

\paragraph{Envelopes from signatures.}
By (R2) and (R5), an edit's effect on a sum is the total of the term changes of the items in its delta, and each term
change has a sign in the corresponding entry of \Cref{tab:fw:terms}, by the term's declared range. For a modified
item, a count term contributes $+1 - 1 = 0$. Since $\mathrm{sgn}(x + y) \in \mathrm{sgn}(x) \boxplus \mathrm{sgn}(y)$,
the sign of the total lies in the sign sum of the entries over the item changes the signature allows. A flagged change
contributes $\{-, 0, +\}$ to every sum its kind writes.

\paragraph{Hashed keys.}
For a signed hashed term whose key the attacker picks from a set $K$, the reachable (bucket, sign) pairs are exactly
the hash and sign values of the keys in $K$, by the term's definition. Enumerating part of $K$ yields witnesses; once
they cover every pair, the answer is exact. With a few dozen buckets, a key that lands in a chosen bucket with a
chosen sign is found by brute force.

\paragraph{Patterns under extension.}
Let a feature count the strings that match a regular expression $P$ that uses no construct looking past the end of
the match (end anchors, word boundaries after the match, lookahead). Whether $P$ matches a substring then depends
only on that substring, which extending the string on the right does not change, so such an extension can only add
matches. The same holds on the left for constructs that look before the match. A pattern that begins with a word
boundary, for example, can gain a match when a string is extended on the right and lose one when it is prefixed with
a letter.

\paragraph{The readout rules (\Cref{tab:fw:readouts}).}
Identity and pools pass the envelope through. A nondecreasing presence readout of one sum changes with a sign in
$E_s \cup \{0\}$, since a count can rise while the bit stays set; of several sums, it can only rise if every sum can
only rise, and only fall if every sum can only fall. A distinct count can only rise if keys are only added, and only
fall if they are only removed. For a ratio $a/b$ with $a \ge 0$ and $b > 0$ before and after, write
$a' = a + \alpha$ and $b' = b + \beta$; then
\[
  \frac{a'}{b'} - \frac{a}{b} = \frac{\alpha\, b - a\, \beta}{b\, b'}.
\]
The denominator is positive, $\alpha b$ has the sign of $\alpha$, which lies in $E_a$, and $-a\beta$ is $0$ if
$a = 0$ and otherwise has a sign in $-E_b$. A normalized vector is a ratio of a bin to the total, and a rate is a ratio
whose denominator, the difference of a max pool and a min pool, changes with a sign in $E_{\max} \boxplus (-E_{\min})$.
The remaining rules (entropy, standard deviation, field transforms, selections, residual features) answer $\{0\}$ only
when nothing they read can change, and every sign otherwise, which is trivially sound.

\paragraph{Proof of Theorem~\ref{thm:fw:directions}.}
For one operation instance on an input of $F$, each sum's change has a sign in its envelope, and by the readout rules
the feature's change has a sign in the rule's output, which lies in the union over operations. Reading the union as a
direction is therefore sound for single instances. If $F$ is closed under $\mathcal{M}$, every intermediate input of a
sequence of edits stays in $F$, the single-instance statement applies at every step, and ``does not fall'' and ``does
not rise'' are transitive; induction on the length of the sequence completes the proof. Without closure, a step can
leave $F$, after which the envelopes say nothing about the next step.

\paragraph{Proof of Corollary~\ref{cor:fw:robustness}.}
By Theorem~\ref{thm:fw:directions}, the feature vectors of $z$ and of any $z' \in \mathrm{Reach}_\mathcal{M}(z)$ differ only
in the allowed directions. The score reads no feature marked both and does not fall in the allowed directions, so
$h(\varphi(z')) \ge h(\varphi(z))$. This is Gimenez et al.'s Proposition~1~\cite{gimenez2025certifiably}, with the
preorder ``$z'$ is reachable from $z$'' and a feature map made monotone by $D$.
